\section{Week 11}
\sol{11.1}{Name the input errors $\Delta x=x-a$ and $\Delta y=y-b$, where the symbol $\Delta$ is simply a label for a difference. Subtraction and regrouping give
\[
 (2x-y)-(2a-b)=2(x-a)-(y-b)=2\Delta x-\Delta y.
\]
Now use the triangle inequality and the absolute-value scaling rule:
\begin{align*}
 |2\Delta x-\Delta y|
 &\le|2\Delta x|+|{-\Delta y}|\\
 &=2|\Delta x|+|\Delta y|\\
 &\le2(0.02)+0.03=0.07.
\end{align*}
The signs of the input errors were not specified, so a guaranteed estimate must allow the worst permitted combination. For $\Delta x=0.02$ and $\Delta y=-0.03$, the output error is $0.04+0.03=0.07$, attaining the bound. For $\Delta x=0.02$ and $\Delta y=0.03$, it is $0.04-0.03=0.01$. In this expression, same-sign errors can partially cancel, while opposite-sign errors reinforce because the second error is subtracted.}
\sol{11.2}{For $a_n=3/(n+1)$ with $n\ge0$, the error from zero is $3/(n+1)$. Given $\eps>0$, choose a positive integer $N>3/\eps$. Then $n\ge N$ implies $n+1>N$ and hence $|a_n-0|=3/(n+1)<3/N<\eps$. The same $N$ works for every later index, which is required for convergence. This choice is sufficient, even though it need not be the smallest possible threshold.}
\sol{11.3}{For positive coordinate dimension $d$, let $d_\infty(x,y)=\max_{1\le i\le d}|x_i-y_i|$. It is nonnegative. It equals zero exactly when every coordinate difference equals zero, hence exactly when $x=y$. Symmetry follows from $|x_i-y_i|=|y_i-x_i|$. For the triangle inequality, each coordinate satisfies
\[
 |x_i-z_i|\le|x_i-y_i|+|y_i-z_i|
 \le d_\infty(x,y)+d_\infty(y,z).
\]
Taking the maximum over $i$ preserves this common bound. Thus it is a metric.}
\sol{11.4}{For any $m>n$, insert the consecutive iterates between the endpoints. Repeated triangle inequalities give
\[
 d(x_m,x_n)\le d(x_m,x_{m-1})+\cdots+d(x_{n+1},x_n)
 \le\sum_{k=n}^{m-1}2(1/3)^k.
\]
The geometric-sum bound from the chapter yields
\[
 \sum_{k=n}^{m-1}2(1/3)^k
 \le\frac{2(1/3)^n}{1-1/3}
 =3(1/3)^n=3^{1-n}.
\]
This bound does not depend on $m$, so it controls all later endpoints together. For any positive tolerance, choose $N$ so $3^{1-N}$ is smaller than that tolerance. When both indices are at least $N$, put the smaller one in the role of $n$; if they are equal their distance is zero. This proves the Cauchy condition. Completeness would be an additional requirement to conclude convergence inside the chosen space.}
\sol{11.5}{The sequence $x_n=2+2^{-(n+1)}$ lies in $(2,3)$ for all $n\ge0$. Its real limit is two because its distance from two is $2^{-(n+1)}\to0$. A convergent real sequence is Cauchy: two sufficiently late terms are each within $\eps/2$ of the limit and hence within $\eps$ of each other. If this sequence converged to a point of $(2,3)$, uniqueness of real limits would force that point to be two, which is outside the interval. Thus this is a Cauchy sequence without a limit in the space.}
\sol{11.6}{One is an upper bound for $(0,1)$. No $b<1$ is an upper bound: if $b<0$, the point $1/2$ exceeds it; if $0\le b<1$, the point $(b+1)/2$ lies in $(0,1)$ and exceeds $b$. Therefore one is the least upper bound. It is not a maximum because it is not a member. Indeed every $x\in(0,1)$ is exceeded by $(x+1)/2$, another member of the interval.}
